%%% SECTION 6: ADAPTED 21 CFR PART 312 FOR PHYSICAL AI IND %%%

The Physical AI Adaption: 21 CFR Part 312 \cite{cfr312-adapt} provides a
comprehensive framework for Investigational New Drug applications in the context
of Physical AI oncology trials, building on the existing regulation
\cite{cfr-part-312}. This is the most extensive of the three regulatory
adaptations, spanning ten subparts (A through J) including a new Subpart J
dedicated entirely to Physical AI systems. The adaptation demonstrates that
Physical AI can be integrated into the IND process without requiring entirely
new legislation, while acknowledging through the creation of Subpart J that
certain aspects of Physical AI are sufficiently novel to warrant their own
regulatory provisions.

\subsection{Scope and Physical AI Definitions}
\label{subsec:cfr312-scope}

Subpart A (General Provisions, \S~312.1-312.10) establishes the scope and
applicability of the adapted regulation. The scope (\S~312.1) extends the
original Part 312 to cover clinical investigations of new drugs that
incorporate Physical AI systems in drug administration, monitoring, or trial
conduct. Applicability (\S~312.2) clarifies which types of Physical AI
involvement trigger IND requirements.

Section 312.3 provides 22 Physical AI definitions, the most comprehensive set
across all three adapted standards. Key definitions include: \textit{Physical
AI Clinical Trial} (a clinical investigation that uses Physical AI systems for
drug administration, patient monitoring, data collection, or trial
coordination); \textit{Phase 0 Simulation Validation} (a mandatory
pre-enrollment phase where Physical AI systems demonstrate safety and
preliminary efficacy through simulation); \textit{Physical AI IND Content} (the
Physical AI-specific information required in an IND submission);
\textit{Cybersecurity Incident} (any unauthorized access, disruption, or
modification of Physical AI systems that may affect patient safety or data
integrity); and \textit{Physical AI Decommissioning} (the process of safely
removing Physical AI systems from clinical trial operations).

Additional provisions in Subpart A address Physical AI labeling (\S~312.6),
requiring clear identification of Physical AI involvement in drug labeling
during investigation; Physical AI promotion restrictions (\S~312.7),
prohibiting promotion of investigational Physical AI drug delivery capabilities;
Physical AI charging provisions (\S~312.8), governing charges for
investigational drugs administered by Physical AI systems; and waiver
procedures (\S~312.10), establishing conditions under which specific Physical
AI requirements may be waived.

\subsection{IND Requirements and Physical AI Content}
\label{subsec:cfr312-ind}

\subsubsection{IND Requirement for Physical AI Trials}

Section 312.20 establishes when an IND is required for Physical AI trials. An
IND is required whenever Physical AI systems are integral to drug
administration (such as robotic IV drug delivery or automated dosing) or when
Physical AI monitoring systems generate data used for safety or efficacy
determinations. The determination of IND requirement considers the Physical AI
system's role in the drug's mechanism of action, the degree to which Physical
AI affects drug exposure or distribution, and whether Physical AI decisions
directly influence dosing or treatment modifications.

\subsubsection{Phases of Investigation}

Section 312.21 adapts the clinical trial phases to include a Phase 0 simulation
validation stage. This is a significant innovation: Physical AI trials include
a mandatory simulation phase before any human subjects are enrolled. Phase 0
requires demonstration that Physical AI systems can safely perform their
intended functions through validated simulations, with performance metrics
meeting established thresholds.

\begin{table}[H]
\centering
\caption{Adapted Clinical Trial Phases for Physical AI}
\label{tab:trial-phases}
\small
\begin{tabularx}{\textwidth}{l X X}
\toprule
\textbf{Phase} & \textbf{Traditional Objective} & \textbf{Physical AI Addition} \\
\midrule
Phase 0 & Not applicable & Simulation validation of Physical AI systems \\
Phase 1 & Safety, dosing & Physical AI system safety with human subjects \\
Phase 2 & Efficacy, side effects & Physical AI-assisted treatment efficacy \\
Phase 3 & Large-scale confirmation & Multi-site Physical AI deployment \\
Phase 4 & Post-approval monitoring & Ongoing Physical AI performance tracking \\
\bottomrule
\end{tabularx}
\end{table}

The Phase 0 requirement is directly supported by the evidence from both A
Cancer Patient's Journey, Physical AI \cite{patient-journey-paper}
(single-patient simulation across ten stages) and the 24-hour simulation from
the Clinical Trial Site Complete Documentation \cite{site-docs} (168 patients,
29 robots, 99.7 percent uptime). These completed simulations demonstrate that
Physical AI systems can be thoroughly validated before real patient contact,
providing a model for how Phase 0 validation should be conducted.

\subsubsection{IND Content and Format}

Section 312.23 adds the Physical AI System Description requirement for IND
submissions. This must include: a complete inventory of all Physical AI systems
to be used, with manufacturer, model, software version, and USL score
\cite{usl-standard}; a description of how each Physical AI system interacts
with the investigational drug (administration, monitoring, safety); validation
evidence including Phase 0 simulation results; a safety analysis addressing
Physical AI-specific risks; and cybersecurity documentation describing
protections for Physical AI systems against unauthorized access or manipulation.

\subsubsection{Protocol Amendments}

Section 312.30 establishes Physical AI amendment triggers: changes to Physical
AI systems that require protocol amendments include replacement of a robot with
a different model or manufacturer, software updates that alter robot behavior
or decision-making, changes to the degree of autonomy (from supervised to
autonomous or vice versa), addition of new Physical AI system types not in the
original protocol, and modifications to safety monitoring parameters.

\subsection{Safety Reporting for Physical AI}
\label{subsec:cfr312-safety}

\subsubsection{Information Amendments}

Section 312.31 on information amendments requires notification to FDA when
Physical AI system updates occur that do not constitute protocol amendments but
that change system characteristics, performance, or safety profile. This
includes routine software patches, firmware updates, calibration adjustments,
and hardware component replacements that fall within pre-approved specifications.

\subsubsection{IND Safety Reporting}

Section 312.32 establishes Physical AI adverse event reporting and cybersecurity
incident reporting. The adaptation creates specific reporting categories for
Physical AI events: Category A covers Physical AI system malfunctions that
directly affect patient safety (reporting within 7 calendar days for serious
events, 15 days for non-serious); Category B covers AI decision errors that
result in incorrect treatment delivery or monitoring failure (reporting within
15 calendar days); Category C covers sensor failures that compromise data
integrity or safety monitoring (reporting within 15 calendar days); and
Category D covers cybersecurity incidents that affect or may affect Physical AI
system integrity (reporting within 24 hours for active breaches, 7 days for
detected vulnerabilities).

\subsubsection{Annual Reports}

Section 312.33 requires that annual reports include a Physical AI performance
summary covering: system uptime and availability statistics, total Physical AI
adverse events by category, software updates and version changes, robot
maintenance and calibration records, cybersecurity assessment results, and
trends in Physical AI system performance over the reporting period.

\subsubsection{IND Withdrawal}

Section 312.38 establishes Physical AI system decommissioning procedures when
an IND is withdrawn. These procedures require safe shutdown of all Physical AI
systems, secure archival of all Physical AI-generated data, notification to
participants about cessation of Physical AI interactions, and physical removal
or secure storage of robot hardware.

\subsection{Administrative Actions for Physical AI Trials}
\label{subsec:cfr312-admin}

Subpart C addresses general administration of Physical AI trials. Section 312.40
establishes general requirements including Physical AI system registration.
Section 312.44 adds Physical AI-specific grounds for IND termination including
repeated Physical AI adverse events demonstrating systemic safety issues,
failure to maintain Physical AI system validation status, and cybersecurity
compromise that cannot be remediated while maintaining patient safety.

Section 312.45 addresses inactive status and Physical AI dormancy, establishing
procedures for placing Physical AI systems in dormant mode during trial pauses
and reactivation procedures ensuring systems meet current safety standards
before resuming operations. Section 312.47 on meetings includes Physical AI
system review as a standing agenda item for pre-IND and end-of-phase meetings.
Section 312.48 provides dispute resolution procedures for Physical AI technical
disputes between sponsors and FDA.

\subsection{Sponsor and Investigator Responsibilities}
\label{subsec:cfr312-responsibilities}

Subpart D (\S~312.50-312.70) adapts sponsor and investigator responsibilities
to Physical AI contexts. Sponsors must ensure Physical AI system procurement
from qualified manufacturers (\S~312.50), maintain Physical AI system
validation status throughout the trial (\S~312.53), and select investigators
with Physical AI competencies (\S~312.55). Investigators must ensure proper
Physical AI system operation (\S~312.60), report Physical AI adverse events
promptly (\S~312.64), and maintain Physical AI system records (\S~312.68).
These responsibilities complement the requirements in the adapted ICH E6(R3)
\cite{ich-adapt} by providing the specific regulatory framework within which
those responsibilities operate.

\subsection{Life-Threatening Illness Provisions}
\label{subsec:cfr312-life-threatening}

Subpart E (\S~312.80-312.88) addresses accelerated approval and expanded access
provisions for Physical AI oncology treatments for life-threatening cancers.
Many cancer patients face conditions where Physical AI may offer faster or
better treatment options, making these provisions particularly relevant. The
adaptation ensures that Physical AI components do not create barriers to
accelerated access when safety data support their use.

\subsection{Expanded Access for Physical AI}
\label{subsec:cfr312-expanded-access}

Subpart I (\S~312.300-312.320) adapts expanded access provisions for Physical
AI, establishing conditions under which Physical AI systems can be made
available outside of clinical trials when sufficient evidence supports their
safety and effectiveness. Expanded access for Physical AI requires
demonstration that the Physical AI system has been validated through Phase 0
simulation, that sufficient clinical trial data support its safety, that no
comparable therapy is available without Physical AI assistance, and that the
requesting site meets minimum PSL requirements.

\subsection{Subpart J: Physical AI Systems}
\label{subsec:cfr312-subpart-j}

Subpart J (\S~312.400-312.405) is entirely new and dedicated to Physical AI
systems. This subpart establishes Physical AI-specific regulatory requirements
that do not fit within the existing subpart structure.

Section 312.400 establishes Physical AI system registration requirements,
requiring all Physical AI systems used in IND trials to be registered in a
centralized database maintained by FDA. Section 312.401 establishes Physical AI
validation standards, including minimum simulation hours, benchmark performance
thresholds, and validation renewal requirements. Section 312.402 establishes
cybersecurity requirements specific to Physical AI systems in clinical trials,
including encryption standards, access control requirements, incident response
procedures, and penetration testing schedules.

Section 312.403 establishes Physical AI interoperability requirements, requiring
systems to support standardized communication protocols (including MCP
\cite{mcp-protocol}) and data formats (including FHIR \cite{hl7-fhir} and
DICOM \cite{dicom}). Section 312.404 establishes Physical AI training
requirements for all personnel who operate, maintain, or oversee Physical AI
systems in clinical trials. Section 312.405 establishes Physical AI
decommissioning standards, including data preservation requirements, hardware
disposal procedures, and participant notification obligations.

The creation of Subpart J demonstrates the balanced approach of this platform:
the existing CFR Part 312 framework can accommodate most Physical AI
requirements through targeted adaptation of existing subparts, but certain
aspects are sufficiently novel to warrant their own provisions. This approach
maintains regulatory continuity with established practice while acknowledging
the unique characteristics of Physical AI systems.
